\chapter{TỔNG KẾT}
\label{ch:tong_ket}

\section{Tóm lược kết quả đạt được}

Đề tài đã xây dựng một quy trình hoàn chỉnh cho bài toán phát hiện biển báo giao thông, từ khâu
kiểm tra dữ liệu thô cho tới bảng số liệu so sánh cuối cùng, và tiến hành bốn nhóm thực nghiệm
trên bộ dữ liệu \textit{pkdarabi/cardetection} gồm 4.969 ảnh với 6.012 hộp bao thuộc 15 lớp.

Đối chiếu với ba câu hỏi nghiên cứu đặt ra ở Chương~\ref{ch:gioi_thieu}:

\textbf{Câu hỏi 1.1} về so sánh giữa hai họ kiến trúc đã được trả lời rõ ràng. YOLOv8n đạt mAP@0,5
bằng 0,9539 và mAP@0,5:0,95 bằng 0,8058 trên tập kiểm tra, với 6,25~MB trọng số và độ trễ 8,83~ms
mỗi ảnh. DETR-ResNet50 chỉ đạt 0,1220 và 0,0988, đồng thời nặng hơn 26,6 lần và chậm hơn 7,2 lần.
Tuy nhiên, kết luận đúng đắn không phải là ``kiến trúc transformer kém hơn'' mà là ``ở ngân sách
dữ liệu và tính toán của đề tài này, DETR chưa được huấn luyện đủ''. Năm yếu tố gây nhiễu đã được
phân tách ở Mục~\ref{sec:vi_sao_detr}, trong đó việc nhánh DETR không hề có tăng cường dữ liệu
trong khi nhánh YOLO có đầy đủ là yếu tố ít được chú ý nhưng nhiều khả năng nặng ký nhất.

\textbf{Câu hỏi 1.2} về kiến trúc nhẹ và nén mô hình cho câu trả lời phủ định ở phần lớn các nhánh.
Cả ba kiến trúc thay thế đều kém hơn mốc tham chiếu của cùng giai đoạn, trong đó YOLO11n tốt nhất
nhưng vẫn thấp hơn 3,81 điểm phần trăm. Ở nhóm nén mô hình, chỉ FP16 là lựa chọn không có nhược điểm: giữ nguyên độ
chính xác trong khi giảm khoảng một nửa dung lượng. Lượng tử hoá INT8 giảm cả độ chính xác lẫn tốc
độ trên môi trường thử nghiệm. Chưng cất tri thức cho kết quả âm rõ rệt.

\textbf{Câu hỏi 1.3} về khoảng cách giữa chỉ số và thực tế là câu hỏi thu được câu trả lời có giá
trị nhất. Phân tích thành phần bộ dữ liệu cho thấy 38,68\% số ảnh là ảnh cắt cận cảnh kiểu GTSRB
với cạnh biển chiếm 66,8\% cạnh ảnh, trong khi 12,98\% là ảnh camera mắt cá với 99,6\% vật thể
thuộc nhóm nhỏ. Phân bố kích thước vì vậy bị lưỡng cực, thiếu hẳn vùng trung gian tương ứng với
biển báo ở khoảng cách lái xe thực tế. Con số 0,9539 do đó phản ánh năng lực trên phân bố của tập
kiểm tra chứ không phản ánh năng lực trên cảnh đường phố thật.

\section{Đóng góp}

Ba đóng góp chính có thể rút ra.

Thứ nhất, về mặt kỹ thuật, đề tài cung cấp một quy trình tái lập được với bộ sinh số ngẫu nhiên
được gieo hạt tập trung, đường dẫn quản lý ở một nơi duy nhất, và bộ kiểm thử tự động gồm 43
trường hợp chạy trên CPU trong khoảng một giây mà không cần GPU hay tải mô hình.

Thứ hai, về mặt thực nghiệm, bảng so sánh mười cấu hình trên ba trục độ chính xác, tốc độ và dung
lượng cho phép người đọc thấy được sự đánh đổi thay vì chỉ một con số duy nhất.

Thứ ba, và quan trọng nhất, đề tài đóng góp một phân tích chỉ ra rằng chỉ số mAP cao trên tập kiểm
tra có thể che giấu hoàn toàn một lỗ hổng năng lực, khi bộ dữ liệu được ghép từ nhiều nguồn có mục
đích khác nhau. Việc kiểm tra phân bố kích thước vật thể \emph{theo từng nguồn} thay vì chỉ trên
toàn tập nên được coi là bước bắt buộc trước khi diễn giải kết quả.

\section{Hạn chế}

Đề tài có những hạn chế sau, được nêu đầy đủ để người đọc đánh giá đúng phạm vi giá trị của kết
quả.

\textbf{Hai nhánh dùng hai cách tính mAP khác nhau.} Nhánh YOLO dùng bộ đánh giá của Ultralytics,
nhánh DETR dùng thư viện torchmetrics. Hai cài đặt có thể khác nhau ở cách nội suy đường cong PR
và cách xử lý ngưỡng IoU, nên chênh lệch giữa hai con số nên được đọc theo xu hướng chứ không theo
từng chữ số. Cách khắc phục triệt để là đưa cả hai về cùng một bộ đánh giá COCO.

\textbf{Bộ dữ liệu có rò rỉ giữa các tập.} Khảo sát ghi nhận 65 trong số 638 ảnh kiểm tra có bản
sao trong tập huấn luyện, tức 10,2\%. Con số 0,9539 vì vậy nên được hiểu là cận trên.

\textbf{Dung lượng mô hình được đo theo hai định nghĩa khác nhau.} Với YOLO, dung lượng lấy theo
kích thước tệp trọng số trên đĩa, vốn được lưu ở độ chính xác 16 bit. Với DETR, dung lượng tính
bằng tổng số byte của tham số ở độ chính xác 32 bit. Tỉ số 26,6 lần vì vậy phóng đại khoảng cách
thực về số lượng tham số, vốn chỉ khoảng 13 lần.

\textbf{Kết quả đến từ hai giai đoạn trên hai cấu hình phần cứng khác nhau.} Nhóm E0 thuộc giai
đoạn 1, các nhóm E1, E2 và E3 thuộc giai đoạn 2. Cấu hình YOLOv8n xuất hiện ở cả hai và cho hai bộ số
khác nhau (0,9539 và 113,2 khung hình mỗi giây so với 0,9633 và 65,3). Báo cáo đã xử lý bằng cách
gắn nhãn giai đoạn cho từng dòng và chỉ so sánh trong nội bộ từng giai đoạn, nhưng đây vẫn là một hạn chế
thực sự: nó khiến không thể đặt DETR và các kiến trúc ở nhóm E1 lên cùng một thước đo. Cách khắc
phục là chạy lại toàn bộ mười cấu hình trong một lượt duy nhất trên cùng một máy.

\textbf{Nhánh DETR không ghi Precision và Recall.} Mô-đun đánh giá DETR chỉ xuất mAP, kích thước mô
hình và tốc độ, nên Bảng so sánh E0 để trống hai ô này. Đây là thiếu sót của quy trình chứ không
phải giới hạn của phương pháp, và sửa được bằng vài dòng mã.

\textbf{Phép đo tốc độ chưa hoàn toàn công bằng.} Nhánh YOLO được đo qua hàm dự đoán đầy đủ, bao
gồm tiền xử lý và NMS, trong khi nhánh DETR chỉ đo lượt truyền xuôi thuần tuý trên đầu vào đã nằm
sẵn trên thiết bị. Chiều thiên lệch này có lợi cho DETR, nên kết luận rằng YOLO nhanh hơn vẫn giữ
nguyên giá trị, thậm chí là một ước lượng dè dặt.

\textbf{Tập lớp bị giới hạn.} Mười lăm lớp chỉ gồm đèn tín hiệu và biển giới hạn tốc độ. Các nhóm
biển cấm, biển nguy hiểm và biển chỉ dẫn hoàn toàn không có trong tập nhãn, nên mô hình không thể
phát hiện chúng dù chất lượng huấn luyện có tốt đến đâu.

\section{Hướng phát triển}

Các hướng phát triển được sắp theo thứ tự ưu tiên, từ chi phí thấp tới chi phí cao.

\textbf{Ưu tiên 1 --- Thống nhất bộ đánh giá cho toàn bộ thực nghiệm.} Nhóm E3 đã dùng pycocotools
và nhờ đó có $AP$ tách theo ba nhóm kích thước --- chỉ số cho thấy mọi khác biệt giữa các cấu hình
đều nằm ở vật thể nhỏ. Áp dụng cùng bộ đánh giá đó cho các nhóm E0, E1 và E2 sẽ vừa bổ sung chỉ số
này cho toàn bộ mười cấu hình, vừa khắc phục hạn chế về hai cách tính mAP khác nhau giữa nhánh YOLO
và nhánh DETR. Đây là việc rẻ nhất trong danh sách vì mã đánh giá đã có sẵn.

\textbf{Ưu tiên 2 --- Làm sạch và cân bằng lại bộ dữ liệu.} Loại bỏ 65 ảnh kiểm tra bị trùng, sau
đó thử huấn luyện lại khi đã loại hoặc giảm mạnh nhóm 1.922 ảnh cắt kiểu GTSRB. Giả thuyết là
mAP tổng sẽ giảm nhưng $AP$ trên nhóm vật thể nhỏ và hiệu năng trên video thật sẽ tăng --- bản thân
việc kiểm chứng giả thuyết này đã là một thí nghiệm có giá trị.

\textbf{Ưu tiên 3 --- Điều chỉnh kiến trúc cho vật thể nhỏ.} Huấn luyện ở kích thước ảnh 960 hoặc
1280 thay vì 640, và bổ sung đầu dự đoán P2 hoạt động ở bước sải 4 thay vì 8. Đây là biện pháp rẻ
hơn nhiều so với chuyển sang một mô hình lớn hơn.

\textbf{Ưu tiên 4 --- Tách bài toán thành hai giai đoạn.} Một bộ phát hiện không phân biệt lớp để
định vị biển báo, kết hợp một bộ phân lớp riêng chạy trên ảnh đã cắt. Cách tổ chức này giải quyết
trực tiếp hiện tượng nhầm lẫn giữa các biển giới hạn tốc độ đã quan sát ở ma trận nhầm lẫn, đồng
thời tận dụng đúng mục đích của 1.922 ảnh cắt kiểu GTSRB.

\textbf{Ưu tiên 5 --- Bổ sung dữ liệu đúng miền.} Không kỹ thuật tăng cường nào thay thế được dữ
liệu thật đúng miền. Gán nhãn vài trăm khung hình từ camera hành trình trong điều kiện giao thông
mục tiêu rồi tinh chỉnh sẽ có tác động lớn hơn mọi biện pháp nêu trên. Quy trình tiết kiệm nhất là
dùng mô hình hiện tại sinh nhãn giả ở ngưỡng tin cậy thấp, sửa tay, rồi tinh chỉnh. Nếu cần bổ
sung dữ liệu công khai, TT100K~\cite{zhu2016tt100k} và Mapillary Traffic
Sign~\cite{ertler2020mapillary} là hai lựa chọn phù hợp vì đều là ảnh cảnh đường thật.

\textbf{Ưu tiên 6 --- Huấn luyện DETR trong điều kiện công bằng.} Tăng số chu kỳ lên gấp mười tới
ba mươi lần, tăng kích thước lô hiệu dụng bằng cách tích luỹ gradient, bổ sung tăng cường dữ liệu
tương đương nhánh YOLO, và cân nhắc dùng Deformable DETR~\cite{zhu2021deformabledetr} vốn hội tụ
nhanh hơn đáng kể. Chỉ khi đó mới có thể kết luận điều gì về so sánh giữa hai họ kiến trúc.
